% !TEX root = ../main.tex
\begin{ouverture}[mouvement]{13}{Décrire un mouvement}{Partie 2 · Mouvement et interactions}
{Position, vitesse, accélération\enspace·\enspace Mouvements rectilignes\enspace·\enspace Repère de Frenet\enspace·\enspace Mouvement circulaire}
\objectif{Définir les vecteurs vitesse et accélération comme des dérivées par rapport au temps.}
\objectif{Établir les coordonnées des vecteurs vitesse et accélération à partir de celles du vecteur position.}
\objectif{Exploiter les coordonnées des vecteurs vitesse et accélération dans le repère de Frenet.}
\objectif{Caractériser le vecteur accélération des mouvements rectilignes et circulaires.}
\objectif{Exploiter une chronophotographie ou une vidéo ; représenter des vecteurs accélération en Python.}
\prerequis{Référentiel, vecteur position, vecteur vitesse (classes de seconde et de première).}
\prerequis{Dérivée d'une fonction ; dérivée de $t\mapsto t^2$, de $\cos$ et de $\sin$.}
\prerequis{Lien entre variation du vecteur vitesse et forces (classe de première).}
\end{ouverture}
\essentiel
\essSection{1}{Position, vitesse et accélération}
\begin{definition}[Vecteurs position, vitesse et accélération]
Dans un référentiel muni d'un repère $(O\,;\vec{\imath},\vec{\jmath},\vec{k})$, la position d'un point M est repérée par le
vecteur $\vect{OM}=x\,\vec{\imath}+y\,\vec{\jmath}+z\,\vec{k}$. Le \textbf{vecteur vitesse} est la dérivée du vecteur position
par rapport au temps, le \textbf{vecteur accélération} la dérivée du vecteur vitesse :
\[\vec{v}=\ddt{\vect{OM}},\qquad \vec{a}=\ddt{\vec{v}}=\frac{\mathrm{d}^2\vect{OM}}{\mathrm{d}t^2}.\]
En coordonnées cartésiennes : $v_x=\ddt{x}$, $v_y=\ddt{y}$, $v_z=\ddt{z}$ et $a_x=\ddt{v_x}$, $a_y=\ddt{v_y}$, $a_z=\ddt{v_z}$. La
vitesse s'exprime en \unit{m.s^{-1}}, l'accélération en \unit{m.s^{-2}}.
\end{definition}
\begin{exemple}
Une balle lancée a pour coordonnées $x(t)=5{,}0\,t$ et $y(t)=-4{,}9\,t^2+8{,}7\,t$ (en mètres et secondes). En dérivant :
$v_x=5{,}0$ et $v_y=-9{,}8\,t+8{,}7$, puis $a_x=0$ et $a_y=-9{,}8$. L'accélération est constante, verticale, vers le bas.
À $t=\qty{0,50}{s}$ : $\vec{v}$ a pour coordonnées $(5{,}0\,;\,3{,}8)$ et pour norme
$v=\sqrt{5{,}0^2+3{,}8^2}=\qty{6,3}{m.s^{-1}}$.
\end{exemple}
\begin{savoirfaire}[Exploiter une chronophotographie]
Les positions $M_0, M_1, M_2\,\dots$ sont enregistrées à intervalles réguliers $\tau$. On approche la vitesse en $M_i$ par
$\vec{v}_i\approx\dfrac{\vect{M_{i-1}M_{i+1}}}{2\tau}$, puis l'accélération par
$\vec{a}_i\approx\dfrac{\vec{v}_{i+1}-\vec{v}_{i-1}}{2\tau}$ : on construit le vecteur $\Delta\vec{v}=\vec{v}_{i+1}-\vec{v}_{i-1}$
au point $M_i$, puis on le divise par $2\tau$.
\end{savoirfaire}
\schema{13-chrono}
\begin{python}
import matplotlib.pyplot as plt
tau = 0.040                                            # s
x = [0.000, 0.048, 0.096, 0.144, 0.192, 0.240, 0.288]  # m
y = [0.000, 0.068, 0.121, 0.157, 0.179, 0.184, 0.174]  # m
vx = [(x[i+1] - x[i-1]) / (2*tau) for i in range(1, 6)]
vy = [(y[i+1] - y[i-1]) / (2*tau) for i in range(1, 6)]
ax = [(vx[i+1] - vx[i-1]) / (2*tau) for i in range(1, 4)]
ay = [(vy[i+1] - vy[i-1]) / (2*tau) for i in range(1, 4)]
print([round(a, 1) for a in ay])
plt.plot(x, y, "o")
plt.quiver(x[2:5], y[2:5], ax, ay)       # accélérations en M2, M3, M4
plt.show()
\end{python}
\begin{console}
[-9.8, -9.7, -9.8]
\end{console}
\noindent\capacite{Capacité numérique}Les listes contiennent les coordonnées d'une balle pointées sur une vidéo. Le
programme calcule les vitesses par la méthode précédente, puis les accélérations, et représente les vecteurs
accélération avec \code{quiver}. Aux erreurs de pointage près, $a_y\approx\qty{-9,8}{m.s^{-2}}$ et $a_x\approx0$ :
l'accélération est verticale, dirigée vers le bas.
\essSection{2}{Mouvements rectilignes}
\begin{propriete}[Accélération d'un mouvement rectiligne]
Le mouvement est rectiligne si la trajectoire est une droite : $\vec{v}$ et $\vec{a}$ sont alors colinéaires à cette
droite.
\begin{itemize}[leftmargin=1.4em,itemsep=0.1em]
\item \textbf{Rectiligne uniforme} : $\vec{v}$ est constant et $\vec{a}=\vec{0}$.
\item \textbf{Rectiligne uniformément varié} : $\vec{a}$ est constant. Sur un axe $(Ox)$ :
$a_x=a$, $v_x(t)=a\,t+v_0$ et $x(t)=\frac12\,a\,t^2+v_0\,t+x_0$.
\item Le mouvement est \textbf{accéléré} si $\vec{a}$ et $\vec{v}$ ont le même sens, \textbf{ralenti} s'ils sont de sens
contraires.
\end{itemize}
\end{propriete}
\begin{exemple}
Une voiture roulant à $v_0=\qty{25}{m.s^{-1}}$ (90~km/h) freine avec une accélération constante de norme
\qty{8,0}{m.s^{-2}}. Sur un axe orienté dans le sens du mouvement, $v_x(t)=25-8{,}0\,t$ s'annule à $t=\qty{3,1}{s}$ ;
la distance de freinage vaut alors $x=25\times3{,}125-\frac12\times8{,}0\times3{,}125^2=\qty{39}{m}$.
\end{exemple}
\begin{attention}
Pour le physicien, « accélérer » signifie que le vecteur vitesse varie, en norme \emph{ou} en direction. Une voiture qui
freine a une accélération, opposée à sa vitesse ; une voiture qui prend un virage à vitesse constante a aussi une
accélération, dirigée vers l'intérieur du virage.
\end{attention}
\essSection{3}{Mouvement circulaire et repère de Frenet}
\begin{definition}[Repère de Frenet]
Pour un point M en mouvement sur un cercle de centre $O$ et de rayon $R$, le repère de Frenet $(M\,;\ut,\un)$ est mobile :
$\ut$ est tangent à la trajectoire, orienté dans le sens du mouvement, et $\un$ est dirigé vers le centre du cercle.
Dans ce repère :
\[\vec{v}=v\,\ut,\qquad \vec{a}=\ddt{v}\,\ut+\frac{v^2}{R}\,\un.\]
La coordonnée tangentielle $a_\mathrm{t}=\ddt{v}$ traduit la variation de la valeur de la vitesse ; la coordonnée normale
$a_\mathrm{n}=\frac{v^2}{R}$, toujours positive, traduit le changement de direction.
\end{definition}
\schema{13-frenet}
\begin{propriete}[Mouvement circulaire uniforme]
Si la valeur $v$ de la vitesse est constante, $a_\mathrm{t}=0$ : l'accélération est \textbf{centripète}, de norme
$a=\frac{v^2}{R}$. Si $T$ est la période (durée d'un tour), $v=\frac{2\pi R}{T}$.
\end{propriete}
\begin{exemple}
Le tambour d'une centrifugeuse de laboratoire, de rayon \qty{0,10}{m}, tourne à 3000 tours par minute, soit 50 tours
par seconde : $v=2\pi\times0{,}10\times50=\qty{31}{m.s^{-1}}$ et $a=\frac{v^2}{R}=\qty{9,9e3}{m.s^{-2}}$, soit environ
\num{1000} fois l'intensité de la pesanteur.
\end{exemple}
\begin{propriete}[Récapitulatif]
\par\smallskip
{\small\renewcommand{\arraystretch}{1.3}
\begin{tabularx}{\linewidth}{@{}lY@{}}\toprule
\otTete{Mouvement} & \otTete{Vecteur accélération}\\\midrule
rectiligne uniforme & nul\\
rectiligne uniformément varié & constant, colinéaire à la trajectoire\\
circulaire uniforme & centripète, de norme constante $v^2/R$\\
circulaire non uniforme & $a_\mathrm{t}=\mathrm{d}v/\mathrm{d}t$ non nulle et $a_\mathrm{n}=v^2/R$ : orienté vers l'intérieur du cercle\\\bottomrule
\end{tabularx}}
\end{propriete}
\begin{experience}[Étudier un mouvement par vidéo]
On filme le lancer d'une balle devant une règle étalon, en plaçant la caméra loin et perpendiculairement au plan du
mouvement. Un logiciel de pointage (ou un programme Python) fournit les coordonnées $x$ et $y$ à chaque image ; on
modélise $x(t)$ par une fonction affine et $y(t)$ par un polynôme du second degré, puis on dérive. On trouve
$a_x\approx0$ et $a_y\approx\qty{-9,8}{m.s^{-2}}$ quelle que soit la façon dont la balle est lancée.
\end{experience}
\begin{rediger}
Toujours préciser le référentiel et le repère. Écrire les dérivées une à une : « $v_y=\ddt{y}=-9{,}8\,t+8{,}7$ ». Pour
caractériser un vecteur, donner sa direction, son sens et sa norme, avec l'unité.
\end{rediger}
\exercices

\rubrique{Découverte}
\begin{blocLivre}{ex-13-1}
\begin{exercice}[1]{Dériver pour décrire un mouvement}
Dans un repère $(O\,;\vec{\imath},\vec{\jmath})$ lié au sol, les coordonnées d'un point M sont, en mètres et en secondes :
\[x(t)=3{,}0\,t+1{,}0\qquad\text{et}\qquad y(t)=-2{,}0\,t^2+4{,}0\,t.\]
\begin{questions}
\item Établir les coordonnées du vecteur vitesse, puis celles du vecteur accélération.
\item Caractériser le vecteur accélération. Le mouvement est-il rectiligne ?
\item Calculer la valeur de la vitesse à $t=0$ et à $t=\qty{1,0}{s}$. À quelle date la vitesse est-elle minimale ?
\item Établir l'équation de la trajectoire $y=f(x)$.
\end{questions}
\end{exercice}
\end{blocLivre}
\begin{blocLivre}{ex-13-2}
\begin{exercice}[1]{Trois chronophotographies}
Les chronophotographies A, B et C représentent les positions successives d'un mobile se déplaçant de gauche à droite,
enregistrées toutes les $\tau=\qty{0,20}{s}$.
\schema{13-ex02}
\begin{questions}
\item Préciser la nature du mouvement dans chaque cas.
\item Pour les mouvements B et C, comparer le sens du vecteur accélération à celui du vecteur vitesse.
\item Calculer la vitesse du mobile A.
\item Pour le mobile B, calculer les vitesses aux positions $M_2$ et $M_4$, puis la valeur de l'accélération en $M_3$.
\end{questions}
\end{exercice}
\end{blocLivre}
\begin{blocLivre}{ex-13-3}
\begin{exercice}[1]{L'accélération du physicien}
Pour chacun des mouvements suivants, préciser si le vecteur accélération est nul ou non et, s'il ne l'est pas,
indiquer sa direction et son sens.
\begin{questions}[label=\otLabel{\alph*.},leftmargin=1.9em,itemsep=0pt,parsep=0.15em]
\item Une voiture roule à 90~km/h sur une ligne droite.
\item La même voiture aborde un virage circulaire à 90~km/h.
\item Un cycliste freine sur une route rectiligne.
\item Une bille est lâchée sans vitesse initiale.
\item La Lune tourne autour de la Terre, sur une orbite supposée circulaire, à vitesse constante.
\item Un ascenseur monte à vitesse constante.
\end{questions}
\begin{questions}
\item Répondre pour chacun des six mouvements.
\item Un conducteur affirme : « dans ce virage, je n'accélère pas, je garde 90~km/h ». Commenter du point de vue du
physicien.
\end{questions}
\end{exercice}
\end{blocLivre}
\begin{blocLivre}{ex-13-4}
\begin{exercice}[1]{Les chaises volantes}
Dans un manège de chaises volantes, chaque nacelle décrit un cercle horizontal de rayon $R=\qty{8,0}{m}$ ; en régime
établi, elle fait un tour en 6,0~s.
\begin{questions}
\item Pourquoi le mouvement d'une nacelle est-il circulaire uniforme en régime établi ?
\item Calculer la valeur de la vitesse d'une nacelle.
\item Caractériser le vecteur accélération : direction, sens et norme. Comparer sa norme à $g=\qty{9,8}{m.s^{-2}}$.
\item Représenter, sur un schéma vu de dessus, le repère de Frenet et les vecteurs vitesse et accélération d'une
nacelle.
\end{questions}
\end{exercice}
\end{blocLivre}
\begin{blocLivre}{ex-13-5}
\begin{exercice}[1]{Lire un vecteur accélération}
Une voiture parcourt un virage circulaire de rayon $R=\qty{50}{m}$, dans le sens indiqué. Le schéma représente son
vecteur accélération en trois points A, B et C.
\schema{13-ex05}
\begin{questions}
\item En quel point la valeur de la vitesse est-elle constante ? croissante ? décroissante ? Justifier.
\item Au point B, la vitesse vaut \qty{12}{m.s^{-1}} et $a_\mathrm{t}=\qty{1,5}{m.s^{-2}}$. Calculer $a_\mathrm{n}$, puis la norme
de l'accélération.
\item Que vaudrait $\vec{a}$ si la voiture roulait en ligne droite à vitesse constante ?
\end{questions}
\end{exercice}
\end{blocLivre}
\rubrique{Entraînement}
\begin{blocLivre}{ex-13-6}
\begin{exercice}[2]{Le départ d'un sprinteur}
Un sprinteur accélère de façon uniforme depuis le départ jusqu'à atteindre \qty{11,6}{m.s^{-1}} au bout de 4,0~s, puis
court à vitesse constante jusqu'à l'arrivée du 100~m.
\schema{13-ex06}
\begin{questions}
\item Quelle est la nature du mouvement dans chaque phase ?
\item Calculer l'accélération pendant la première phase et écrire les équations horaires $v(t)$ et $x(t)$.
\item Quelle distance le sprinteur a-t-il parcourue au bout de 4,0~s ? Retrouver ce résultat à l'aide de l'aire sous la
courbe $v(t)$.
\item En déduire le temps réalisé sur 100~m.
\end{questions}
\end{exercice}
\end{blocLivre}
\begin{blocLivre}{ex-13-7}
\begin{exercice}[2]{La distance d'arrêt}
Un automobiliste roule à 90~km/h sur une route rectiligne. Il aperçoit un obstacle : il lui faut 1,0~s pour réagir,
puis il freine avec une accélération constante de norme \qty{7,5}{m.s^{-2}} sur route sèche, ou
\qty{4,0}{m.s^{-2}} sur route mouillée.
\begin{questions}
\item Quelle distance la voiture parcourt-elle pendant le temps de réaction ? Quelle est la nature du mouvement ?
\item Pendant le freinage, on prend l'origine des dates et des positions au début du freinage. Écrire $v(t)$ et $x(t)$.
\item Calculer la durée du freinage et la distance de freinage sur route sèche, puis la distance d'arrêt.
\item Calculer la distance d'arrêt sur route mouillée. Commenter.
\end{questions}
\end{exercice}
\end{blocLivre}
\begin{blocLivre}{ex-13-8}
\begin{exercice}[2]{Une chute filmée}
Une balle est lâchée sans vitesse initiale devant une règle. On pointe sa position sur un axe $(Oy)$ vertical,
\textbf{orienté vers le bas}, toutes les $\tau=\qty{0,050}{s}$ ; le programme calcule les vitesses aux points
intermédiaires.
\begin{python}
tau = 0.050
y = [0.000, 0.012, 0.049, 0.110, 0.196, 0.307, 0.441]   # m
v = [(y[i+1] - y[i-1]) / (2 * tau) for i in range(1, 6)]
print([round(vi, 2) for vi in v])
\end{python}
\begin{console}
[0.49, 0.98, 1.47, 1.97, 2.45]
\end{console}
\begin{questions}
\item \capacite{Capacité numérique}Quelle approximation de la vitesse le programme utilise-t-il ? À quelles dates
correspondent les cinq vitesses affichées ?
\item Tracer l'allure de $v_y$ en fonction du temps. Quelle est la nature du mouvement ?
\item Déterminer la valeur de l'accélération. Commenter.
\item Ajouter au programme les lignes qui calculent et affichent les accélérations.
\end{questions}
\end{exercice}
\end{blocLivre}
\begin{blocLivre}{ex-13-9}
\begin{exercice}[2]{Sur un rond-point}
Une voiture s'engage sur un rond-point circulaire de rayon $R=\qty{20}{m}$. Sa vitesse passe de \qty{5,0}{m.s^{-1}} à
\qty{10,0}{m.s^{-1}} en 4,0~s, la coordonnée tangentielle de l'accélération restant constante.
\begin{questions}
\item Calculer $a_\mathrm{t}$.
\item Calculer $a_\mathrm{n}$ lorsque la vitesse vaut \qty{8,0}{m.s^{-1}}, puis la norme de l'accélération.
\item Représenter, en ce point, les vecteurs $\vec{v}$ et $\vec{a}$ dans le repère de Frenet ; calculer l'angle entre
$\vec{a}$ et $\un$.
\item Comment évolue $a_\mathrm{n}$ pendant la phase d'accélération ?
\end{questions}
\end{exercice}
\end{blocLivre}
\begin{blocLivre}{ex-13-10}
\begin{exercice}[2]{La Station spatiale internationale}
Dans le référentiel géocentrique, la Station spatiale internationale décrit une orbite circulaire à 420~km d'altitude ;
elle fait un tour en 92,8~min.
\begin{donnees}
Rayon de la Terre : $R_\mathrm{T}=\qty{6371}{km}$.
\end{donnees}
\begin{questions}
\item Calculer le rayon de l'orbite, puis la valeur de la vitesse de la station, en \unit{m.s^{-1}} et en km/h.
\item Caractériser son vecteur accélération.
\item Combien de tours la station fait-elle par jour ?
\item Comparer la valeur de l'accélération à celle de la pesanteur au sol. Ce résultat sera interprété au chapitre~16.
\end{questions}
\end{exercice}
\end{blocLivre}
\begin{blocLivre}{ex-13-11}
\begin{exercice}[2]{Représenter l'accélération d'un mouvement circulaire}
Un point M décrit un cercle de rayon $R=\qty{2,0}{m}$ à vitesse constante, avec une période $T=\qty{4,0}{s}$. Le
programme suivant calcule une valeur approchée de son accélération à la date $t=\qty{0,5}{s}$.
\begin{python}
from math import cos, sin, pi
R, T, tau = 2.0, 4.0, 0.1                 # m, s, s
def x(t): return R * cos(2 * pi * t / T)
def y(t): return R * sin(2 * pi * t / T)
t = 0.5
ax = (x(t + tau) - 2 * x(t) + x(t - tau)) / tau**2
ay = (y(t + tau) - 2 * y(t) + y(t - tau)) / tau**2
print(round(ax, 3), round(ay, 3))
\end{python}
\begin{console}
-3.482 -3.482
\end{console}
\begin{questions}
\item \capacite{Capacité numérique}Calculer les coordonnées de M à $t=\qty{0,5}{s}$ et placer le point sur un schéma.
\item Justifier que $\frac{x(t+\tau)-2x(t)+x(t-\tau)}{\tau^2}$ est une valeur approchée de $a_x$.
\item Représenter le vecteur accélération obtenu. Vers quel point est-il dirigé ?
\item Comparer sa norme à la valeur $v^2/R$ attendue.
\end{questions}
\end{exercice}
\end{blocLivre}
\rubrique{Objectif bac}
\begin{blocLivre}{ex-13-12}
\begin{exercice}[3]{La centrifugeuse des pilotes}
\begin{docu}{1}{S'entraîner aux fortes accélérations}
Les pilotes de chasse et les astronautes s'entraînent dans des centrifugeuses : une cabine, fixée au bout d'un bras de
longueur $R=\qty{7,6}{m}$, tourne dans un plan horizontal. L'accélération subie s'exprime en multiples de $g$ ; un
pilote entraîné supporte environ $8\,g$ pendant quelques secondes.
\end{docu}
\begin{docu}{2}{Le protocole de montée en régime}
Partant de l'arrêt, la valeur de la vitesse de la cabine augmente avec une accélération tangentielle constante pendant
20~s, jusqu'à la vitesse correspondant à $8\,g$, puis reste constante.
\end{docu}
\begin{donnees}
$g=\qty{9,81}{m.s^{-2}}$.
\end{donnees}
\begin{questions}
\item Quelle est la nature du mouvement de la cabine en régime établi ? Caractériser son vecteur accélération.
\item Calculer la vitesse de la cabine correspondant à $8\,g$, puis le nombre de tours par minute.
\item Calculer l'accélération tangentielle pendant la montée en régime.
\item \initiative Au milieu de la montée en régime, déterminer la norme de l'accélération de la cabine et l'angle qu'elle
fait avec la direction du bras.
\end{questions}
\end{exercice}
\end{blocLivre}
\begin{blocLivre}{ex-13-13}
\begin{exercice}[3]{Le lancer franc}
\begin{docu}{1}{Une vidéo analysée}
Un basketteur tire un lancer franc. Le pointage vidéo du centre du ballon, dans un repère d'origine au sol, sous le
point de lâcher, avec $(Ox)$ horizontal vers le panier et $(Oy)$ vertical vers le haut, est modélisé par :
\[x(t)=4{,}78\,t\qquad\text{et}\qquad y(t)=-4{,}90\,t^2+5{,}70\,t+2{,}10\quad\text{(m, s)}.\]
\end{docu}
\begin{docu}{2}{Le panier}
Le centre de l'anneau se trouve à $x=\qty{4,60}{m}$, à une hauteur de \qty{3,05}{m}.
\end{docu}
\begin{questions}
\item Établir les coordonnées des vecteurs vitesse et accélération du ballon. Caractériser l'accélération.
\item Calculer la valeur de la vitesse initiale et l'angle de tir par rapport à l'horizontale.
\item À quelle date le ballon atteint-il le sommet de sa trajectoire ? À quelle hauteur ?
\item \initiative Le centre du ballon passe-t-il par le centre de l'anneau ? Justifier.
\end{questions}
\end{exercice}
\end{blocLivre}
\begin{blocLivre}{ex-13-14}
\begin{exercice}[3]{Le looping d'un grand huit}
\begin{docu}{1}{Une boucle en forme de goutte}
Les loopings des montagnes russes ne sont pas circulaires : leur rayon de courbure est grand en bas et petit en haut.
Au point bas A, où le rayon de courbure vaut \qty{15}{m}, le wagon passe à \qty{22}{m.s^{-1}} ; au sommet B, où il vaut
\qty{8,0}{m}, il passe à \qty{12}{m.s^{-1}}.
\end{docu}
\schema{13-ex14}
\begin{docu}{2}{Limite de confort}
Pour le grand public, les constructeurs limitent l'accélération à environ $5\,g$, avec $g=\qty{9,8}{m.s^{-2}}$.
\end{docu}
\begin{questions}
\item En assimilant la trajectoire, au voisinage de A et de B, à des arcs de cercle, calculer la coordonnée normale de
l'accélération en A et en B. Les exprimer en multiples de $g$.
\item En A et en B, la coordonnée tangentielle de l'accélération est nulle. Représenter les vecteurs accélération en ces
deux points.
\item \initiative Que vaudrait l'accélération en A pour une boucle circulaire de rayon \qty{8,0}{m} parcourue à la même
vitesse ? Justifier le choix d'une boucle en forme de goutte.
\end{questions}
\end{exercice}
\end{blocLivre}
\begin{blocLivre}{ex-13-15}
\begin{exercice}[3]{Le TGV dans les courbes}
\begin{docu}{1}{Des courbes à très grand rayon}
Sur une ligne à grande vitesse, où les TGV circulent à 320~km/h, le rayon des courbes est d'au moins 7,0~km. Sur une
ligne classique, certaines courbes ont un rayon de 1,2~km seulement.
\end{docu}
\begin{docu}{2}{Le freinage}
Lancé à 320~km/h, un TGV a besoin d'environ 3,3~km pour s'arrêter.
\end{docu}
\begin{questions}
\item Calculer la coordonnée normale de l'accélération d'un TGV roulant à 320~km/h dans une courbe de 7,0~km de rayon.
\item À quelle vitesse doit-il aborder une courbe de 1,2~km de rayon pour subir la même accélération normale ? Commenter.
\item \initiative En supposant le freinage uniformément ralenti, calculer l'accélération et la durée du freinage.
\item Comparer les deux accélérations calculées à $g=\qty{9,8}{m.s^{-2}}$. Pourquoi les passagers les supportent-ils
sans difficulté ?
\end{questions}
\end{exercice}
\end{blocLivre}
